\section*{第 \ref{ch:g10:habitable-planet} 章——宜居行星}

\begin{solution}{exo:g10:habitable-planet:1}
理想辐射体向太空辐射的功率恰好
等于行星从恒星吸收的功率时的温度。
地球表面更暖（\qty{288}{K}，相比约 \qty{255}{K}）
是因为温室气体吸收了地球发出的部分红外线，并将其中一部分送回地面。
\end{solution}

\begin{solution}{exo:g10:habitable-planet:2}
金星：实测温度 \qty{737}{K}，\omterm{def:g10:habitable-planet:equilibrium-temperature}{平衡温度}约为 \qty{227}{K}，
升温约为 \qty{510}{K}。
\end{solution}

\begin{solution}{exo:g10:habitable-planet:3}
$S = 1361/1.52^2 \approx \qty{589}{W/m^2}$（资料表给出
\qty{586}{W/m^2}，采用精确平均距离）。
\end{solution}

\begin{solution}{exo:g10:habitable-planet:4}
\qty{1}{bar}，\qty{330}{K}：液体（低于沸点，约
\qty{373}{K}）。\qty{0.001}{bar}，\qty{250}{K}：冰（压强
位于冰–蒸气曲线之上，该曲线在 \qty{250}{K} 时约为 \qty{0.0008}{bar}）。
\qty{100}{bar}，\qty{700}{K}：蒸气——温度高于临界温度，
因此不可能存在液体。
\end{solution}

\begin{solution}{exo:g10:habitable-planet:5}
绕太阳以外的恒星运行的行星。凌日法：行星
每次公转时从恒星前方经过，使恒星亮度略有下降。摆动法：行星的
引力使恒星沿视线方向来回移动，
从而周期性地改变其光的颜色。
\end{solution}

\begin{solution}{exo:g10:habitable-planet:6}
$T = (1361 \times 0.706/(4 \times 5.67\times10^{-8}))^{1/4} \approx
\qty{255}{K} \approx \qty{-18}{\degreeCelsius}$。当 $A = 0$ 时：
$\approx \qty{278}{K} \approx \qty{5}{\degreeCelsius}$。
\end{solution}

\begin{solution}{exo:g10:habitable-planet:7}
$\sqrt 4 = 2$：从 $2 \times 0.99 = \qty{1.98}{au}$ 到 $2 \times 1.70 =
\qty{3.40}{au}$。在 \qty{1}{au} 处，行星接收到的辐照度
为地球的四倍，远在内缘之内，温度过高。
\end{solution}

\begin{solution}{exo:g10:habitable-planet:8}
$S = 1361/0.723^2 \approx \qty{2604}{W/m^2}$；
$T_{\mathrm{eq}} = (2604 \times 0.23/(4\sigma))^{1/4} \approx \qty{227}{K}$。
其表面约高出 \qty{510}{K}，原因在于其厚重的
二氧化碳大气（\qty{92}{bar}）产生了巨大的温室效应。
\end{solution}

\begin{solution}{exo:g10:habitable-planet:9}
在 \qty{6.36}{mbar} 的压强下，水仅在约 \qty{0}{\degreeCelsius}
至 \qty{0.5}{\degreeCelsius} 之间为液体。在 \qty{10}{\degreeCelsius} 时，
压强低于沸腾曲线：水会汽化为蒸气；夜间则凝固。
\end{solution}

\begin{solution}{exo:g10:habitable-planet:10}
$11.2/2.4 \approx 4.7$。气体分子要逃离地球，所需速度
几乎比逃离月球时快五倍；经过数十亿年，月球微弱的引力
使其气体逸散，而地球则保留了其中的大部分。
\end{solution}

\begin{solution}{exo:g10:habitable-planet:11}
$T = (1361 \times 0.50/(4\sigma))^{1/4} \approx \qty{234}{K}$，约
\qty{-39}{\degreeCelsius}。冰会反射大量阳光，使行星保持足够寒冷以维持冰层：
这是一个正反馈回路。
\end{solution}

\begin{solution}{exo:g10:habitable-planet:12}
$T_{\mathrm{eq}} = (S_0(1-A)/(4\sigma d^2))^{1/4} = (S_0/(4\sigma))^{1/4}
(1-A)^{1/4} (d^2)^{-1/4}$；$(1361/(4\sigma))^{1/4} \approx \qty{278.3}{K}$，以及
$(d^2)^{1/4} = \sqrt d$。距离为四倍时：$T$ 除以 $\sqrt 4 = 2$。
\end{solution}

\begin{solution}{exo:g10:habitable-planet:13}
$\sqrt d = 278.3 \times 0.706^{1/4}/273 \approx 0.934$，于是 $d \approx
\qty{0.873}{au}$，位于内缘之内（\qty{0.99}{au}）。\omterm{def:g10:habitable-planet:equilibrium-temperature}{平衡
温度}未考虑温室效应：带大气的真实行星
比 $T_{\mathrm{eq}}$ 更暖，因此会过热，并在
更靠外的约 \qty{0.99}{au} 处就已失去水分。
各边界来自整个大气的模型，而非仅由 $T_{\mathrm{eq}}$ 决定。
\end{solution}

\begin{solution}{exo:g10:habitable-planet:14}
液态水存在的温度范围是熔点（约 \qty{0}{\degreeCelsius}）与
沸点（约 \qty{7}{\degreeCelsius}）之间。在 \qty{15}{\degreeCelsius} 时，湖泊
将会沸腾。
\end{solution}

\begin{solution}{exo:g10:habitable-planet:15}
$\sqrt{0.5} \approx 0.707$：从 $0.707 \times 0.99 \approx \qty{0.70}{au}$
到 $0.707 \times 1.70 \approx \qty{1.20}{au}$。该行星恰好位于
外缘处。边界是模型结果，比例换算只是近似，因此
无法确定它位于边界之内还是之外。
\end{solution}

\begin{solution}{pb:g10:habitable-planet:1}
\textbf{1.} $0.00151/0.0485^2 \approx 0.642$。

\textbf{2.} $0.642 \times 1361 \approx \qty{874}{W/m^2}$。

\textbf{3.} 介于两者之间：约为金星的三分之一，
约为火星的 1.5 倍。

\textbf{4.} $365.25/11.18 \approx 32.7$。

\textbf{5.} $(874/(4 \times 5.67\times10^{-8}))^{1/4} \approx \qty{249}{K}$。

\textbf{6.} $(874 \times 0.23/(4\sigma))^{1/4} \approx \qty{173}{K}$。

\textbf{7.} 两者皆非：\qty{249}{K} 仍低于 \qty{273}{K}。亮色的
云层或冰盖会使\omterm{def:g10:habitable-planet:equilibrium-temperature}{平衡温度}下降超过
\qty{75}{K}：行星表面覆盖什么，与其距离同等重要。

\textbf{8.} $\sqrt{0.00151} \approx 0.0389$。

\textbf{9.} $0.99 \times 0.0389 \approx \qty{0.0385}{au}$。

\textbf{10.} $1.70 \times 0.0389 \approx \qty{0.066}{au}$。

\textbf{11.} 是：$0.0385 < 0.0485 < 0.066$。

\textbf{12.} $d_\odot = 0.0485/0.0389 \approx \qty{1.25}{au}$：位于
太阳\omterm{def:g10:habitable-planet:habitable-zone}{宜居带}的内侧，介于地球与火星之间。

\textbf{13.} $1.055^{1/3} \approx 1.018$。

\textbf{14.} $11.2 \times \sqrt{1.055/1.018} \approx \qty{11.4}{km/s}$，约
与地球相当：其引力足以保留大气（是否
拥有大气取决于其演化历史）。

\textbf{15.} 恒星的摆动只揭示其运动中沿
我们视线方向的分量；由于轨道倾角未知，所得的
质量仅在轨道面与视线平行时为真实质量，否则只是下限。

\textbf{16.} 三相点的压强为 \qty{611.657}{Pa}，
约 \qty{0.006}{bar}。

\textbf{17.} $(874 \times 0.706/(4\sigma))^{1/4} \approx \qty{228}{K}$。

\textbf{18.} $288 - 228 = \qty{60}{K}$：比邻星b需要
\textbf{约 \qty{60}{K}}的温室增温，几乎是地球 \qty{33}{K} 的两倍。
\end{solution}
